\documentclass{resume} % Use the custom resume.cls style
\usepackage{hyperref}
\geometry{left=0.75in,top=0.5in,right=0.75in,bottom=0.5in}
\newcommand{\tab}[1]{\hspace{.2667\textwidth}\rlap{#1}}
\newcommand{\itab}[1]{\hspace{0em}\rlap{#1}}
\name{Leyi Zhu}
\address{leyi@nyu.edu}

\begin{document}

%----------------------------------------------------------------------------------------
% EXPERIENCE SECTION
%----------------------------------------------------------------------------------------

\begin{rSection}{Experience}

\begin{rSubsection}{Google}{June 2026 -- Present}{Software Engineer, YouTube Shorts Feed Optimization}{}
\item Work on learned ranking functions for the YouTube Shorts recommendation feed.
\end{rSubsection}

\begin{rSubsection}{Tencent America, Pixel Lab}{June 2022 -- Aug. 2022}{Graphics Research Intern}{}
\item Developed a method for mesh connectivity optimization under user-defined distortion constraints, with tools to transfer geometric data across mesh connectivities.
\item Improved accuracy and convergence of common geometric distortion energies (e.g., AMIPS, Symmetric Dirichlet) via adaptive-precision arithmetic.
\end{rSubsection}

\end{rSection}

%----------------------------------------------------------------------------------------
% EDUCATION SECTION
%----------------------------------------------------------------------------------------

\begin{rSection}{Education}

{\bf New York University} \hfill {\em Sept. 2019 -- Feb. 2026}\\
Ph.D. in Computer Science, specializing in Computer Graphics \hfill \\
Courant Institute of Mathematical Sciences \\
Thesis: {\em Surface Parameterization and Remeshing with Guarantees}

{\bf University of Science and Technology of China (USTC)} \hfill {\em Sept. 2015 -- July 2019}\\
B.S. in Information \& Computational Science \hfill Ranking: 1/46\\
Special Class for the Gifted Young

\end{rSection}

%----------------------------------------------------------------------------------------
% PUBLICATIONS SECTION
%----------------------------------------------------------------------------------------

\begin{rSection}{Selected Publications}

{\bf BijectiveRemesh: Maintaining Bijective Mappings for Data Transfer Across Remeshed Manifolds},
{\bf L. Zhu}, M. Tao, Y. Hu, D. Panozzo, D. Zorin,
\textit{arXiv preprint}, 2026

{\bf Which Cross Fields can be Quadrangulated? Global Parameterization from Prescribed Holonomy Signatures},
H. Shen, {\bf L. Zhu}, R. Capouellez, D. Panozzo, M. Campen, D. Zorin,
\textit{ACM Trans. Graph. (SIGGRAPH)}, 2022

{\bf Efficient and Robust Discrete Conformal Equivalence with Boundary},
M. Campen, R. Capouellez, H. Shen, {\bf L. Zhu}, D. Panozzo, D. Zorin,
\textit{ACM Trans. Graph. (SIGGRAPH Asia)}, 2021

{\bf Simulation and Visualization of Ductile Fracture with the Material Point Method},
S. Wang, M. Ding, T. Gast, {\bf L. Zhu}, S. Gagniere, C. Jiang, J. Teran,
\textit{ACM SIGGRAPH/Eurographics SCA}, 2019 \quad {\bf (Best Paper Award)}

\end{rSection}

%----------------------------------------------------------------------------------------
% HONORS SECTION
%----------------------------------------------------------------------------------------

\begin{rSection}{Honors}

MacCracken Fellowship, New York University \hfill 2019 -- 2025\\
National Scholarship, USTC \hfill 2016, 2018\\
Huawei Scholarship, USTC \hfill 2017

\end{rSection}

%----------------------------------------------------------------------------------------
% SKILLS SECTION
%----------------------------------------------------------------------------------------

\begin{rSection}{Skills}

\begin{tabular}{ @{} >{\bfseries}l @{\hspace{6ex}} l }
Programming Languages & C/C++, Python, MATLAB, Shell \\
Machine Learning & Learning-to-Rank, Recommender Systems, Distributed Training \\
Software \& Tools & PyTorch, Mathematica, Blender, ParaView, \LaTeX \\
\end{tabular}

\end{rSection}

\end{document}
